% ============================================================================
% Simulation section: national-AI-validation
% Requires: amsmath, amssymb, booktabs, graphicx
% All floats use [h!] as requested.
% ============================================================================
\providecommand{\navdir}{national-AI-validation}
\providecommand{\navfigdirR}{\navdir/figures/r_publication}
\providecommand{\navfigdirPy}{\navdir/figures/python_validation}
\providecommand{\navincludegraphics}[2][]{%
  \IfFileExists{\navfigdirR/#2.pdf}%
    {\includegraphics[#1]{\navfigdirR/#2.pdf}}%
    {\includegraphics[#1]{\navfigdirPy/#2.pdf}}%
}

\section{Simulation Study: Probability-Based National Validation of a Fixed Clinical AI System}
\label{sec:sim-national-ai-validation}

\subsection{Purpose and policy question}

This simulation provides a single empirical demonstration of the chapter's
estimation-grounded validation framework. A fixed clinical AI system produces
a binary alert for clinical deterioration within 24 hours. Before national
deployment, a public authority must determine whether sensitivity is at least
\(\tau=0.85\) in the target-country population and whether the same requirement
is met among harder-to-predict patients. The complete predictive pipeline,
including preprocessing, missing-data handling, fitted parameters, operating
threshold, calibration, and post-processing, is locked before the audit is
designed.

The experiment evaluates complete estimation strategies
\begin{equation}
  \eta_k=(d_k,\widehat\theta_k,V_k),
  \label{eq:nav-strategy}
\end{equation}
where \(d_k\) is the probability construction process,
\(\widehat\theta_k\) the estimator, and \(V_k\) the uncertainty procedure.
The principal purpose is to show that the evidentiary value of a realised
Objective Reference Point (ORP) does not arise from its observations alone.
It depends on the target parameter and on whether the estimator and variance
procedure are compatible with the construction process.

This is a target-adequacy experiment. The WEIRD-country development setting
provides motivation only; source performance and the External Transportability
Criterion are not simulated.

\subsection{Fixed target population and estimands}

One finite target population was generated and held fixed throughout the
Monte Carlo experiment. It contains all \(N_T=25{,}585\) eligible admissions
in \(H=400\) hospitals during one common, pre-specified audit period. There is
no calendar-month sampling stage. Hospitals belong to a frame-known design
stratum \(H_h\in\{0,1\}\), with 300 standard hospitals and 100 priority
hospitals. The priority stratum has a higher expected prevalence of difficult
patients and deterioration events. Patient difficulty is represented by the
separate analysis variable \(Q_{hi}\in\{0,1\}\). Both values of \(Q\) occur
in both hospital strata.

For admission \((h,i)\), let \(Y_{hi}\) indicate true deterioration,
\(A_{hi}\) denote the alert produced by the locked system, and
\(T_{hi}=A_{hi}Y_{hi}\) denote a true positive. The national sensitivity
estimand is
\begin{equation}
  \theta_T
  =
  \Pr_T(A=1\mid Y=1)
  =
  \frac{\sum_h\sum_i A_{hi}Y_{hi}}
       {\sum_h\sum_i Y_{hi}}.
  \label{eq:nav-national-target}
\end{equation}
For \(q\in\{0,1\}\), the patient-subgroup estimands are
\begin{equation}
  \theta_{T,q}
  =
  \Pr_T(A=1\mid Y=1,Q=q).
  \label{eq:nav-subgroup-target}
\end{equation}

The finite-population truths were
\[
  \theta_T=0.8784,
  \qquad
  \theta_{T,0}=0.9480,
  \qquad
  \theta_{T,1}=0.7368.
\]
National and easier-subgroup performance therefore exceed the adequacy
threshold, whereas harder-subgroup performance is clearly below it.

\begin{table}[h!]
\centering
\small
\caption{Fixed finite-population target quantities.}
\label{tab:nav-population-truths}
\begin{tabular}{lrrrr}
\toprule
Estimand & Eligible patients & Positive cases & Event prevalence & Sensitivity \\
\midrule
National & 25,585 & 4,819 & 0.188 & 0.8784 \\
Easy & 20,300 & 3,231 & 0.159 & 0.9480 \\
Hard &  5,285 & 1,588 & 0.300 & 0.7368 \\
\bottomrule
\end{tabular}
\par\vspace{2pt}
\begin{minipage}{0.97\linewidth}\footnotesize The harder subgroup is defined at the patient level and occurs in both hospital design strata.\end{minipage}
\end{table}



Figure~\ref{fig:nav-population-structure} shows why the hospital stratum is
useful for prospective enrichment without being identical to the patient
subgroup. Priority hospitals contain more harder patients and more outcome
events, but sensitivity is also lower in that stratum.

\begin{figure}[h!]
  \centering
  \navincludegraphics[width=0.90\linewidth]{nav_fig1_population_structure}
  \caption{Finite-population structure by hospital design stratum.}
  \label{fig:nav-population-structure}
\end{figure}

\subsection{Target, observation, and realised ORP distributions}

The target distribution \(P_T\) is the finite-population empirical
distribution assigning equal mass to every eligible admission. For strategy
\(k\), the design-induced unweighted observation distribution is
\begin{equation}
  P_{\mathcal O,k}(h,i)
  =
  \frac{\pi_{hi}^{(k)}}
       {\sum_g\sum_r \pi_{gr}^{(k)}},
  \label{eq:nav-po}
\end{equation}
where \(\pi_{hi}^{(k)}\) is the first-order probability that admission
\((h,i)\) appears in the ORP. Under hospital sampling followed by a census of
eligible admissions, \(\pi_{hi}^{(k)}=\pi_h^{(k)}\). The associated
unweighted sensitivity parameter is
\begin{equation}
  \theta_{\mathcal O,k}
  =
  \frac{\sum_h\sum_i \pi_{hi}^{(k)}A_{hi}Y_{hi}}
       {\sum_h\sum_i \pi_{hi}^{(k)}Y_{hi}}.
  \label{eq:nav-theta-o}
\end{equation}
This quantity is distinct from both \(\theta_T\) and the empirical estimate
computed from one realised ORP.

In the enriched design, one half of the selected hospitals is drawn from each
hospital stratum. Because the frame contains three times as many standard as
priority hospitals, the priority-hospital inclusion probability is three times
larger. The exact enriched unweighted national parameter is
\(\theta_{\mathcal O,3}=0.8419\), producing the reference--target discrepancy
\begin{equation}
  \operatorname{RTD}_3
  =
  \theta_{\mathcal O,3}-\theta_T
  =-0.0365.
  \label{eq:nav-rtd}
\end{equation}
The corresponding discrepancies are \(-0.0056\) for easier patients and
\(-0.0258\) for harder patients.

\begin{table}[h!]
\centering
\small
\caption{Target and design-induced unweighted observation parameters for the enriched ORP.}
\label{tab:nav-observation-parameters}
\begin{tabular}{lrrr}
\toprule
Estimand & $\theta_T$ & $\theta_{\mathcal O,3}$ & RTD \\
\midrule
National & 0.8784 & 0.8419 & -0.0365 \\
Easy & 0.9480 & 0.9424 & -0.0056 \\
Hard & 0.7368 & 0.7110 & -0.0258 \\
\bottomrule
\end{tabular}
\par\vspace{2pt}
\begin{minipage}{0.97\linewidth}\footnotesize The inclusion-probability ratio between priority and standard hospitals is three for every enriched sample-size condition, so the observation parameters do not depend on $m$.\end{minipage}
\end{table}



\begin{figure}[h!]
  \centering
  \navincludegraphics[width=0.90\linewidth]{nav_fig2_target_observation_parameters}
  \caption{Target parameters and design-induced unweighted observation parameters for the enriched ORP.}
  \label{fig:nav-target-observation}
\end{figure}

\subsection{Design--estimator--variance strategies}

Five strategies were compared. N1 is an ideal patient-level benchmark that
requires a complete patient frame. N2 and N4 are operationally realistic
hospital-frame designs. N3 and N4 use exactly the same enriched realised ORP
within every replication but different point estimators. N4 and N5 use the
same ORP and exactly the same point estimate but different uncertainty
procedures.

For weighted hospital strategies, stratum totals are estimated independently.
For \(j\in\{0,1\}\),
\begin{equation}
  \widehat T_j
  =
  \frac{H_j}{m_j}\sum_{h\in s_j}T_h,
  \qquad
  \widehat Y_j
  =
  \frac{H_j}{m_j}\sum_{h\in s_j}Y_h,
  \label{eq:nav-stratum-totals}
\end{equation}
and national sensitivity is estimated as
\begin{equation}
  \widehat\theta
  =
  \frac{\widehat T_0+\widehat T_1}
       {\widehat Y_0+\widehat Y_1}.
  \label{eq:nav-combined-ratio}
\end{equation}
Subgroup estimators use subgroup-specific numerator and denominator totals in
exactly the same way. No target positive-case shares are supplied to the
estimator.

\begin{table}[h!]
\centering
\small
\caption{Probability construction, estimator, and uncertainty procedure for the five strategies.}
\label{tab:nav-strategies}
\begin{tabular}{lp{0.19\linewidth}p{0.20\linewidth}p{0.23\linewidth}p{0.18\linewidth}}
\toprule
Strategy & Construction process & Point estimator & Variance procedure & Role \\
\midrule
N1 & Patient SRSWOR & Ordinary sensitivity ratio & Patient-SRS ratio linearisation & Ideal self-weighting benchmark \\
N2 & Proportional stratified hospital SRS & Design-weighted combined ratio & Stratified cluster Taylor linearisation & Realistic target-aligned design \\
N3 & Enriched stratified hospital SRS & Unweighted pooled ratio & Cluster-aware variance around $\theta_{\mathcal O,3}$ & Deliberate estimator mismatch \\
N4 & Same enriched ORP as N3 & Design-weighted combined ratio & Stratified cluster Taylor linearisation & Compatible enrichment \\
N5 & Same enriched ORP and point estimate as N4 & Design-weighted combined ratio & Naive patient-level variance & Deliberate uncertainty mismatch \\
\bottomrule
\end{tabular}
\end{table}



The total hospital sample sizes were
\[
  m\in\{40,80,120,160\}.
\]
N2 used proportional allocations
\((30,10)\), \((60,20)\), \((90,30)\), and \((120,40)\). N3--N5 used
enriched allocations \((20,20)\), \((40,40)\), \((60,60)\), and
\((80,80)\). Every selected hospital contributed all eligible admissions
from the common audit period. N1 patient sample sizes were chosen to equal the
expected number of admissions observed under proportional hospital sampling.

N2 and N4 used a stratified hospital-level Taylor linearisation for the ratio.
N3 used a cluster-aware variance around the unweighted parameter it actually
targeted. N5 instead treated weighted positive patients as independent and
unclustered. Logit-transformed 95\% confidence intervals were used throughout.

The PR evaluation tolerances were
\begin{equation}
  \varepsilon_b=0.03,
  \qquad
  \varepsilon_v
  =
  \left(\frac{0.03}{1.96}\right)^2.
  \label{eq:nav-pr-tolerances}
\end{equation}
Bias and precision were assessed from 5,000 repeated samples per design and
sample-size condition. Empirical interval coverage was reported separately,
not folded into the PR variance condition. Because TAC requires a reliable
interval in addition to point-estimation PR, a formal TAC classification was
issued only when the PR conditions, target-interval validity, and declared
variance-procedure compatibility were all satisfied.

\subsection{Results}

\paragraph{Same target parameter under different probability designs.}
N1, N2, and N4 all recovered national sensitivity when paired with compatible
estimators. Across all sample sizes, N1's absolute national bias was below
\(0.0004\), while N4's was below \(0.0003\). N2 was also effectively unbiased.
The designs differed in precision because patient SRS and clustered hospital
sampling have different variance structures, but all three provided a valid
route to the same target parameter once their design information was used.

\begin{table}[h!]
\centering
\small
\caption{Monte Carlo performance for national sensitivity.}
\label{tab:nav-national-performance}
\begin{tabular}{rrlrrrrr}
\toprule
$m$ & Strategy & Bias & ESD & RMSE & Mean SE & Coverage & Evidence \\
\midrule
40 & N1 & +0.0001 & 0.0141 & 0.0141 & 0.0141 & 0.953 & Pass \\
40 & N2 & -0.0000 & 0.0238 & 0.0238 & 0.0230 & 0.922 & Fail \\
40 & N3 & -0.0371 & 0.0254 & 0.0450 & 0.0251 & 0.629 & Fail \\
40 & N4 & -0.0008 & 0.0207 & 0.0207 & 0.0205 & 0.946 & Fail \\
40 & N5 & -0.0008 & 0.0207 & 0.0207 & 0.0144 & 0.828 & Fail \\
80 & N1 & +0.0001 & 0.0093 & 0.0093 & 0.0094 & 0.955 & Pass \\
80 & N2 & -0.0003 & 0.0157 & 0.0157 & 0.0155 & 0.941 & Fail \\
80 & N3 & -0.0367 & 0.0157 & 0.0399 & 0.0157 & 0.288 & Fail \\
80 & N4 & -0.0002 & 0.0135 & 0.0135 & 0.0135 & 0.949 & Pass \\
80 & N5 & -0.0002 & 0.0135 & 0.0135 & 0.0102 & 0.863 & Fail \\
120 & N1 & +0.0000 & 0.0072 & 0.0072 & 0.0072 & 0.949 & Pass \\
120 & N2 & +0.0002 & 0.0121 & 0.0121 & 0.0119 & 0.945 & Pass \\
120 & N3 & -0.0366 & 0.0109 & 0.0382 & 0.0109 & 0.045 & Fail \\
120 & N4 & -0.0000 & 0.0101 & 0.0101 & 0.0100 & 0.947 & Pass \\
120 & N5 & -0.0000 & 0.0101 & 0.0101 & 0.0083 & 0.892 & Fail \\
160 & N1 & +0.0000 & 0.0058 & 0.0058 & 0.0058 & 0.951 & Pass \\
160 & N2 & +0.0001 & 0.0096 & 0.0096 & 0.0096 & 0.944 & Pass \\
160 & N3 & -0.0365 & 0.0073 & 0.0373 & 0.0073 & 0.000 & Fail \\
160 & N4 & -0.0001 & 0.0079 & 0.0079 & 0.0078 & 0.946 & Pass \\
160 & N5 & -0.0001 & 0.0079 & 0.0079 & 0.0072 & 0.925 & Fail \\
\bottomrule
\end{tabular}
\par\vspace{2pt}
\begin{minipage}{0.97\linewidth}\footnotesize Evidence is marked Pass only when the bias and precision tolerances are met, target interval coverage is acceptable, and the declared variance procedure is design-compatible.\end{minipage}
\end{table}



\begin{figure}[h!]
  \centering
  \navincludegraphics[width=0.98\linewidth]{nav_fig3_bias}
  \caption{Bias by strategy, hospital sample size, and sensitivity estimand. Dotted horizontal lines mark the pre-specified bias tolerance.}
  \label{fig:nav-bias}
\end{figure}

\paragraph{Same enriched ORP, different point estimators.}
N3 and N4 contained exactly the same selected hospitals and admissions in
every replication. Nevertheless, N3's national bias remained approximately
\(-0.0365\) at every sample size, closely matching the exact RTD. Its empirical
standard deviation decreased from \(0.0259\) at \(m=40\) to \(0.0073\) at
\(m=160\), but its mean estimate remained centred near
\(\theta_{\mathcal O,3}=0.8419\). N4 used the same ORP with inverse-inclusion
weights and remained centred on \(\theta_T=0.8784\).

This comparison directly separates the realised ORP from the estimation
strategy. Weighting did not alter the observed units or transform the ORP into
a different dataset. It altered the point estimator's inferential target.

\begin{figure}[h!]
  \centering
  \navincludegraphics[width=0.90\linewidth]{nav_fig9_large_wrong_orp}
  \caption{Increasing sample size concentrates N3 around the design-induced observation parameter rather than the national target.}
  \label{fig:nav-large-wrong-orp}
\end{figure}

\paragraph{Same ORP and point estimate, different uncertainty.}
N4 and N5 had numerically identical estimates in every replication. Their
uncertainty statements differed substantially. At \(m=80\), the empirical
standard deviation of the common national estimator was \(0.0135\). N4's mean
estimated standard error was also \(0.0135\), with 95\% coverage of \(0.947\).
N5's mean standard error was only \(0.0102\), and coverage fell to \(0.861\).
At \(m=120\), coverage was \(0.942\) for N4 and \(0.890\) for N5.

The narrower N5 intervals also produced more mechanically decisive TAC
results. For example, at \(m=80\), the mechanical adequate-decision frequency
was 68.6\% under N5 compared with 49.1\% under N4. These additional decisions
were not additional evidence; they resulted from ignoring the hospital
clustering that generated the ORP.

\begin{table}[h!]
\centering
\small
\caption{National sensitivity: same enriched samples and point estimates, different uncertainty procedures.}
\label{tab:nav-variance-comparison}
\begin{tabular}{rrlrrrr}
\toprule
$m$ & Strategy & ESD & Mean SE & SE/ESD & Coverage & Mean width \\
\midrule
40 & N4 & 0.0207 & 0.0205 & 0.991 & 0.946 & 0.0812 \\
40 & N5 & 0.0207 & 0.0144 & 0.693 & 0.828 & 0.0566 \\
80 & N4 & 0.0135 & 0.0135 & 1.000 & 0.949 & 0.0530 \\
80 & N5 & 0.0135 & 0.0102 & 0.754 & 0.863 & 0.0399 \\
120 & N4 & 0.0101 & 0.0100 & 0.999 & 0.947 & 0.0394 \\
120 & N5 & 0.0101 & 0.0083 & 0.825 & 0.892 & 0.0325 \\
160 & N4 & 0.0079 & 0.0078 & 0.987 & 0.946 & 0.0306 \\
160 & N5 & 0.0079 & 0.0072 & 0.911 & 0.925 & 0.0282 \\
\bottomrule
\end{tabular}
\par\vspace{2pt}
\begin{minipage}{0.97\linewidth}\footnotesize N4 and N5 have numerically identical point estimates in every replication. Differences arise only from the uncertainty procedure.\end{minipage}
\end{table}



\begin{figure}[h!]
  \centering
  \navincludegraphics[width=0.90\linewidth]{nav_fig5_coverage_n4_n5}
  \caption{Empirical target-parameter coverage for the same enriched point estimator under compatible and naive variance procedures.}
  \label{fig:nav-coverage}
\end{figure}

\paragraph{Enrichment and harder-subgroup evidence.}
The enriched strategy increased harder-subgroup positive cases by
approximately 49\% relative to proportional hospital sampling at every sample
size. At \(m=160\), the mean count increased from 635 under N2 to 947 under
N4. The empirical standard deviation of harder-subgroup sensitivity decreased
from \(0.0222\) to \(0.0147\), a reduction of about 34\%. N4 consequently
passed the pre-specified hard-subgroup precision criterion at \(m=160\), while
N2 remained evidentially insufficient for that estimand.

\begin{table}[h!]
\centering
\small
\caption{Effect of enrichment on harder-subgroup evidence.}
\label{tab:nav-hard-evidence}
\begin{tabular}{rrlrrrr}
\toprule
$m$ & Strategy & Mean hard positives & ESD & RMSE & Precision & Evidence \\
\midrule
40 & N2 & 159.0 & 0.0545 & 0.0545 & Fail & Fail \\
40 & N4 & 236.5 & 0.0433 & 0.0434 & Fail & Fail \\
80 & N2 & 318.1 & 0.0358 & 0.0358 & Fail & Fail \\
80 & N4 & 473.5 & 0.0273 & 0.0273 & Fail & Fail \\
120 & N2 & 475.9 & 0.0276 & 0.0276 & Fail & Fail \\
120 & N4 & 710.0 & 0.0201 & 0.0201 & Fail & Fail \\
160 & N2 & 635.0 & 0.0222 & 0.0222 & Fail & Fail \\
160 & N4 & 946.0 & 0.0146 & 0.0145 & Pass & Pass \\
\bottomrule
\end{tabular}
\par\vspace{2pt}
\begin{minipage}{0.97\linewidth}\footnotesize At $m=160$, N4 passes the pre-specified precision criterion for harder-subgroup sensitivity, whereas N2 remains evidentially insufficient.\end{minipage}
\end{table}



\begin{figure}[h!]
  \centering
  \navincludegraphics[width=0.90\linewidth]{nav_fig6_hard_positive_yield}
  \caption{Harder-subgroup positive-case yield under proportional and enriched hospital sampling.}
  \label{fig:nav-hard-yield}
\end{figure}

\paragraph{PR was estimand-relative.}
The same strategy did not receive a single global classification. N4 supported
national sensitivity from \(m=80\) onward but supported harder-subgroup
sensitivity only at \(m=160\). N2 supported national sensitivity at
\(m=120\) and \(m=160\) but did not achieve the required harder-subgroup
precision. N3 failed the national bias criterion at every sample size even
when its sampling variance became small. These results operationalise the
claim that PR is relative to a parameter and tolerance, not a descriptive
property of the ORP.

\begin{table}[h!]
\centering
\small
\caption{Predictive Representativity and uncertainty diagnostics at $m=160$.}
\label{tab:nav-pr-status}
\begin{tabular}{llrrrrr}
\toprule
Strategy & Estimand & Bias & Precision & Coverage & $V_k$ compatible & Full evidence \\
\midrule
N1 & Easy & Pass & Pass & Pass & Pass & Pass \\
N1 & Hard & Pass & Pass & Pass & Pass & Pass \\
N1 & National & Pass & Pass & Pass & Pass & Pass \\
N2 & Easy & Pass & Pass & Pass & Pass & Pass \\
N2 & Hard & Pass & Fail & Pass & Pass & Fail \\
N2 & National & Pass & Pass & Pass & Pass & Pass \\
N3 & Easy & Pass & Pass & Fail & Pass & Fail \\
N3 & Hard & Pass & Pass & Fail & Pass & Fail \\
N3 & National & Fail & Pass & Fail & Pass & Fail \\
N4 & Easy & Pass & Pass & Pass & Pass & Pass \\
N4 & Hard & Pass & Pass & Pass & Pass & Pass \\
N4 & National & Pass & Pass & Pass & Pass & Pass \\
N5 & Easy & Pass & Pass & Pass & Fail & Fail \\
N5 & Hard & Pass & Pass & Pass & Fail & Fail \\
N5 & National & Pass & Pass & Fail & Fail & Fail \\
\bottomrule
\end{tabular}
\end{table}



\begin{figure}[h!]
  \centering
  \navincludegraphics[width=0.92\linewidth]{nav_fig7_evidential_status}
  \caption{Full evidential status by strategy, sample size, and estimand.}
  \label{fig:nav-evidential-status}
\end{figure}

\paragraph{TAC followed PR and uncertainty validation.}
At \(m=160\), N4 produced a formal national TAC decision of Adequate in
92.3\% of replications and Inconclusive in 7.7\%. For harder-subgroup
sensitivity, where the truth was 0.7368, N4 produced Not adequate in 100\% of
replications. N2 remained evidentially insufficient for the harder-subgroup
TAC because its precision criterion failed. N3 was not assigned a national TAC
classification because it estimated the wrong national parameter, and N5 was
not assigned a formal classification because its declared uncertainty
procedure was incompatible with the clustered design.

\begin{table}[h!]
\centering
\small
\caption{Formal TAC decision frequencies at $m=160$.}
\label{tab:nav-tac}
\begin{tabular}{llrrrr}
\toprule
Strategy & Estimand & Adequate & Inconclusive & Not adequate & Evidentially insufficient \\
\midrule
N1 & Easy & 1.000 & 0.000 & 0.000 & 0.000 \\
N1 & Hard & 0.000 & 0.000 & 1.000 & 0.000 \\
N1 & National & 0.997 & 0.003 & 0.000 & 0.000 \\
N2 & Easy & 1.000 & 0.000 & 0.000 & 0.000 \\
N2 & Hard & 0.000 & 0.000 & 0.000 & 1.000 \\
N2 & National & 0.782 & 0.218 & 0.000 & 0.000 \\
N3 & Easy & 0.000 & 0.000 & 0.000 & 1.000 \\
N3 & Hard & 0.000 & 0.000 & 0.000 & 1.000 \\
N3 & National & 0.000 & 0.000 & 0.000 & 1.000 \\
N4 & Easy & 1.000 & 0.000 & 0.000 & 0.000 \\
N4 & Hard & 0.000 & 0.000 & 1.000 & 0.000 \\
N4 & National & 0.911 & 0.089 & 0.000 & 0.000 \\
N5 & Easy & 0.000 & 0.000 & 0.000 & 1.000 \\
N5 & Hard & 0.000 & 0.000 & 0.000 & 1.000 \\
N5 & National & 0.000 & 0.000 & 0.000 & 1.000 \\
\bottomrule
\end{tabular}
\par\vspace{2pt}
\begin{minipage}{0.97\linewidth}\footnotesize Mechanical threshold results are suppressed whenever the full evidential gate fails.\end{minipage}
\end{table}



\begin{figure}[h!]
  \centering
  \navincludegraphics[width=0.96\linewidth]{nav_fig8_tac}
  \caption{Formal TAC decision frequencies at \(m=160\), with evidential insufficiency shown separately from interval-based inconclusiveness.}
  \label{fig:nav-tac}
\end{figure}

\subsection{Interpretation}

The numerical results support six conclusions. First, \(P_T\) defines where
the performance claim applies, whereas \(P_{\mathcal O,k}\) describes the
unweighted composition induced by the design. Second, the empirical ORP is a
realised reference object rather than a target parameter. Third, the same ORP
can support different inferential conclusions when paired with different
estimators, as demonstrated by N3 and N4. Fourth, the same point estimator can
have different evidentiary reliability under different variance procedures,
as demonstrated by N4 and N5. Fifth, prospective enrichment can improve
harder-subgroup precision while preserving national inference when inclusion
probabilities are known and the estimator is compatible. Sixth, increasing
sample size reduces sampling error but does not remove a fixed
reference--target discrepancy.

The practical implication is that national validation should be planned from
the policy question backward. The target population and performance parameter
must be specified first; the hospital frame and allocation should then be
chosen to support both national and subgroup questions; and the estimator and
variance procedure must be declared as part of the strategy. PR is the
evidential gatekeeper. TAC is the subsequent decision rule and cannot repair a
failure of target alignment or uncertainty validity.

The conclusions concern sensitivity at one locked operating threshold. They
do not establish global model safety, calibration, discrimination, or clinical
utility, each of which would require its own estimand-compatible validation
strategy.
